% 本文件由 LaTeX 排版 Agent 根据有效 translation.json 自动生成。
% author=Clement of Alexandria; work_id=0211; title=残篇（亚历山大城之克莱孟）

\chapter{残篇（亚历山大城之克莱孟）}

% structure: p0001 source=h1 target=omitted reason=page-title-already-rendered-by-parent

% structure: p0002 source=h2 target=section reason=relative-heading-rank; base=h2

\section{一、——摘自卡西奥多鲁斯的拉丁文译本。}

\Emph{一、——论伯多禄前书注释}。\par

\Emph{伯前一3} \Quote{愿我们的主耶稣基督的天主和父受赞美，祂因自己的大仁慈，重生了我们。}因为如果天主由物质生了我们，祂后来藉着生活的进步，重生了我们。\par

\Quote{我们主之父，藉着耶稣基督的复活：}那位按照你们的信德，在我们内复活；另一方面，因我们不信的行为，祂在我们内死亡。因为祂又说过，灵魂在此生中决不再返回肉身；那已成为天使性的，不会成为不义或邪恶，以致没有藉着披上肉身再次犯罪的机会；但在复活中，灵魂返回肉身，两者按其本性互相结合，藉着各自的构成，像石头建筑一样和谐一致。\par

此外，伯多禄说，\BibleRef{伯前 2:5} \Quote{你们也像活石，被建成一座属灵的殿宇；}意指天上守护的天使居所。\Quote{因为你们，}他说，\Quote{藉天主的德能、信德和默观，被保守着，要接受你们信德的终向，即灵魂的救恩。}\par

由此可见，灵魂并非本性不朽，而是藉天主的恩宠，通过信仰、正义和知识成为不朽。\Quote{关于这救恩，}他说，\BibleRef{伯前 1:10}\Quote{那些先知们已经寻求和探究过，}等等。由此宣告，先知以智慧发言，基督之神在他们内，按照对基督的拥有，并从属于基督。因为天主藉着总领天使和同类的天使工作，他们被称为基督之神。\par

\Quote{这些现在}他说，\BibleRef{伯前 1:12}\Quote{传报给你们，那传福音给你们的人。}先知所行、多数人未觉察的旧事，如今藉着传福音者启示给你们。\Quote{因为给你们，}他说，\Quote{藉着被派遣的圣神，显明了这些事；}那就是护慰者，主曾论及祂说：\Quote{我若不去，祂就不会来。}\BibleRef{若 16:7}\Quote{连天使们，}经上说，\Quote{也渴望窥探这些事；}并非如多数人所猜疑的背叛的天使，而是一个神圣的真理：天使们渴望获得那完美的益处。\par

\Quote{凭借宝贵的血，}他说，\BibleRef{伯前 1:19}\Quote{如同无瑕疵、无玷污的羔羊的血。}此处他论及古代的肋未和司祭庆典，但意指一个藉着正义而纯洁、献给天主的灵魂。\par

\Quote{确实在创世以前就被预知。}因为祂在万物之先被预知，因为祂是基督。\Quote{但在这末世才显现出来}藉着肉身的诞生。\Quote{你们重生，不是由于可朽坏的种子。}因此，与身体一同产生的灵魂是可朽坏的，如有些人所想。\par

\Quote{但主的话永远长存：}既是预言，也是神圣教训。\par

\Quote{但是你们是特选的种族，王家的司祭。}\BibleRef{伯前 2:9}我们藉天主的拣选成为特选的种族，这是十分清楚的。祂说王家，因为我们是蒙召来统御，并属于基督；司祭则是因为藉祈祷和训导所作的祭献，藉此赢得献给天主的灵魂。\par

\Quote{祂被辱骂，}他说，\Quote{却不还口；受虐待，却不恐吓。}主以祂的良善和忍耐如此行事。\Quote{却把自己交托给那按不义审判祂的：}或指祂自己，如此看自己，这是一种换位。祂确实把自己交给那些按不义法律审判的人，因为祂对他们无用，因为祂是义人；或者，祂把那些不义审判、无端坚持处死祂的人交托给天主，使他们因受惩罚而受教导。\par

\Quote{凡愿意爱慕生命，并看见美好日子的人；}\BibleRef{伯后 3:10}即愿意成为永恒不朽的人。祂称主为生命，称那日子为美好，即神圣。\par

\Quote{因为上主的眼看顾义人，祂的耳垂听他们的祈祷：}祂意指圣神多方面的眷顾。\Quote{上主的颜面却敌视作恶的人：}即审判、报复或显明。\par

\Quote{但是要在你们心里尊主基督为圣，}他说。因为你们在主祷文中这样说：\Quote{愿你的名被尊为圣。}\BibleRef{玛 6:9}\par

\Quote{因为基督也曾一次为我们的罪受了苦，义人代替不义的人，为要把我们引到天主面前；就肉身而论，祂被处死；就神性而言，祂复活了。}祂说这些话，是将他们引向信仰。这就是说，祂在我们的灵里成为活的。\par

\Quote{祂来临，}他说，\Quote{也曾向那些曾经不信的人宣讲。}他们没有看见祂的形状，却听见了祂的声音。\par

\Quote{当时天主的容忍}在持续。天主如此良善，以致藉着救恩的教导产生效果。\par

\Quote{藉着耶稣基督的复活：}即那藉着信德在我们内成就的复活。\par

\Quote{天使服从于祂，}为第一等级；\Quote{率领者}服从，为第二等级；\Quote{德能者}也服从，据称属于第三等级。\par

\Quote{谁将要交账，}他说，\BibleRef{伯前 4:5}\Quote{给那位准备审判活人死人的。}\par

这些人藉着先前的审判而受到训练。因此，祂加上说：\BibleQuote{为此，福音也曾传报给死人}——也就是说，传报给我们这些从前不信的人。\BibleQuote{为使他们在肉身方面，按人受审判，}他说，\BibleQuote{却在神魂方面，按天主而生活。}这是因为他们曾经背离了信仰；当他们仍在肉身中时，他们按先前的审判受审判，为使他们悔改。因此，他也加上说：\BibleQuote{为使他们按天主在神魂内生活。}保禄也是如此；因为当他说\BibleQuote{我已将这人交给撒殚，为使他的神魂得以得救}时，他也说了类似的话；\BibleRef{格前 5:5}也就是说，\BibleQuote{如同天主各种恩宠的好管家。}同样，保禄也说：\BibleQuote{天主在古时，曾多次并以多种方式，藉着先知对我们的祖先说过话。}\BibleRef{希 1:1}\par

\BibleQuote{你们欢喜吧，}经上说，\BibleQuote{因为你们有分于基督的苦难。}也就是说，如果你们是义人，你们为义受苦，正如基督为义受苦一样。\BibleQuote{你们是有福的，因为圣神，即祂荣耀与德能的圣神，安息在你们身上。}这里的物主代词\Quote{祂的}也指示一个属天使的神：即那些藉着信德与义德，使天主得荣耀的人，就是天主的光荣，按照被引进的圣徒的进展，归于可敬的荣耀。\Quote{圣神在我们身上}可如此理解：即那藉着信德来到灵魂上，如同心灵的优雅和灵魂之美。\par

\BibleQuote{因为时候到了，审判要从天主的家开始。}因为审判将在指定的迫害中临到这些人。\par

\BibleQuote{但那赐万般恩宠的天主，}他说。\BibleRef{伯前 5:10}\BibleQuote{万般恩宠，}他说，因为祂是善的，且是一切美善的施予者。\par

\Quote{我的儿子马尔谷问候你们。}马尔谷，伯多禄的门徒，当伯多禄在罗马公开宣讲福音，面对凯撒的一些骑士，并引证许多关于基督的见证时，为使他们能记住所说的话，便完整地写下了所谓的\WorkTitle{马尔谷福音}。同样，路加也可从文体辨认出，他编写了\WorkTitle{宗徒大事录}，并翻译了\WorkTitle{保禄致希伯来人书}。\par

% structure: p0028 source=h2 target=section reason=relative-heading-rank; base=h2

\section{二、——论\WorkTitle{犹达书}的注释}

犹达，这位写了公函的人，是若瑟众子的兄弟，非常虔诚，虽然知道主的近亲关系，却没有称自己为祂的兄弟。但他说了什么？\BibleQuote{犹达，耶稣基督的仆人}——视祂为主；却\BibleQuote{雅各伯的兄弟。}这是真的；祂是祂的兄弟，（若瑟之子）。\BibleQuote{因为有些人已经偷偷进来，不虔敬的人，他们从起初就被预定并注定要受我们天主的审判。}并非他们要成为不虔敬的人，而是他们既然不虔敬，就被注定受审判。\BibleQuote{因为主天主，}他说，\BibleQuote{曾一次把百姓从埃及领出来，后来却毁灭了那些不信的人。}也就是说，为藉惩罚来训练他们。因为他们确实受了惩罚，并为了那些得救的人而灭亡，直到他们回头归向主。\BibleQuote{至于那些天使，}他说，\BibleQuote{没有保持自己尊位的}——即他们藉着提升所得到的——\BibleQuote{却离弃了自己的居所}，意指那居所，即天和星辰，他们便成为并被称为背弃者。\BibleQuote{祂用锁链，在幽暗中拘留了他们，以待大日的审判。}祂指靠近地的地方，即黑暗的空气。他称\Quote{锁链}为他们所站立的尊荣的丧失，以及对软弱之物的贪欲；因为他们被自己的贪欲所束缚，无法悔改。\BibleQuote{如同索多玛和哈摩辣，}他说……主以此表明宽恕已赐予；且他们在受管教后已悔改。\BibleQuote{同样，这些做梦的人，}他说——即那些在想象中梦见贪欲和邪恶欲望的人，将不是真正为善、且超越一切善的事物视为善——\BibleQuote{污秽肉身，藐视主权，并亵渎尊荣者}，即那唯一的主，真正是我们的主耶稣基督，唯独堪受赞美。他们\Quote{亵渎尊荣者}，即亵渎天使。\par

\BibleQuote{当总领天使弥额尔，与魔鬼争辩，辩论关于梅瑟的身体时。}这里他确认了梅瑟的被提升天。这里被称为弥额尔的，是藉着一位与我们相近的天使与魔鬼争辩者。\par

\BibleQuote{但这些，}他说，\BibleQuote{他们不知道的事，他们便亵渎；而他们本性所知道的事，如同无理性的牲畜，就在这些事上败坏自己。}他意指他们吃、喝、沉溺于不洁，并说他们做其他与动物共有的、缺乏理性的事。\par

\BibleQuote{他们有祸了！}他说，\BibleQuote{因为他们走了加音的道路。}同样，我们也因罪的相似而落在亚当的罪下。\BibleQuote{云彩，}他说，\BibleQuote{没有水，他们自身不拥有神圣而有成果的话语。}因此，他说，\BibleQuote{这类人被风与猛烈的气流所吹动。}\BibleQuote{秋天的树，}他说，\BibleQuote{没有果实，}——即不信者，不结信德的果实。\BibleQuote{两次死亡，}他说：一次，当他们因过犯而犯罪时；第二次，当按天主预定的审判交付惩罚时；因为即使各人不立即承受，也该算为死亡。\BibleQuote{狂浪，}他说，\BibleQuote{的汹涌大海。}这些话他意指外邦人的生命，其结局是可憎的野心。\BibleQuote{游荡的星辰，}——即，他意指那些犯错和背教者，属于从天使座位坠落的那类星辰——\BibleQuote{为他们，}由于他们的背教，\BibleQuote{永远的黑暗幽冥被保留。}他又说，\BibleQuote{亚当的第七代子孙\Person{哈诺客}，也曾预言这些。}这些话他证实了预言。\par

\BibleQuote{那些，}他说，\BibleQuote{分别出来}——从信者与不信者分开，按他们自己的不信被定罪。又说那些从肉身分开的。他说，\BibleQuote{没有灵性的人；}即，因信德而来的灵，藉着实践义德而来的灵。\par

\BibleQuote{但是你们，可爱的，}他说，\BibleQuote{要在你们至圣的信德上，在圣神内，建筑自己。}\BibleQuote{有些人，}他说，\BibleQuote{你们要拯救，把他们从火中拉出来；}\BibleQuote{对有些人却要怀着恐惧怜悯，}即，教导那些陷入火中的人如何自救。\BibleQuote{憎恨，}他说，\BibleQuote{那被玷污的、属肉体的衣裳：}即灵魂的衣裳；被玷污的衣裳是被情欲污染的灵。\par

\BibleQuote{愿荣耀归于那能保守你们不失脚，并且在祂荣耀显现时，使你们无瑕无疵地喜乐站在祂面前的那一位。}祂的荣耀显现：他意指在众天使面前，被呈现为无瑕无疵，成了天使。当\Person{达尼尔}论及人民并来到主面前时，他并没有这样说，因为他看见了天主：因为凡心不纯洁的人不可能看见天主；但他这样说，是因为人民所做的一切都在天主眼中，并且是显露给祂的；也就是说，没有什么是向主隐藏的。\par

如今，在据\WorkTitle{马尔谷福音}中，主被大司祭问及祂是否是基督，那受赞颂的天主之子，祂回答说：\BibleQuote{我是；你们将看见人子坐在大能者的右边。}但大能者指圣天使。再者，当祂说\BibleQuote{在天主的右边，}祂指同一类受造物，因着天使和神圣的大能者同等的相似，他们被称为天主之名。因此，祂说祂坐在右边；即，祂安息于卓越的尊荣中。然而，在其他福音中，据说祂没有答复大司祭询问祂是否为天主之子。但祂说了什么？\BibleQuote{你说得是。}回答得恰好。因为若祂说：正如你所理解的，祂就说了不真实的话，没有承认自己是天主子；因为他们并不持守这个看法；但说\BibleQuote{你说得是}，祂说了实话。因为他们所不了解、却用言语表达的事，祂就承认是真实的。\par

% structure: p0037 source=h2 target=section reason=relative-heading-rank; base=h2

\section{三、——对\WorkTitle{若望一书}的注释。}

\BibleRef{若一 1:1}\BibleQuote{那从起初就有的，我们亲眼所见，亲耳所闻。}\par

继\WorkTitle{若望福音}之后，并与之一致，这封书信也包含了属灵的原则。\par

因此，他所说的\BibleQuote{从起初}，长老解释如下：受生的开始并不与创造者的开始分离。因为当他说\BibleQuote{那从起初就有的}，他触及了圣子的无始生，祂与圣父共同存在。所以，有一个\OriginalTerm{圣言}{Word}蕴含着无始的永恒；正如圣言自身，即天主子，由于本质同等，与圣父为一，是永恒且非受造的。祂永远是圣言，由\BibleQuote{在起初已有圣言}这句话所表明。但藉着\BibleQuote{我们亲眼所见}，他意指主在肉身中的临在，\BibleQuote{我们的手也摸过}，他说，\BibleQuote{生命之圣言。}他不仅意指祂的肉身，也指圣子的德能，如同穿透最深处之地的阳光——这阳光在肉身中到来，成为门徒可触摸的。因此，在传统中记载，\Person{若望}触摸那外在的肉身本身，将手深探入其中，而肉身的坚实并未构成阻碍，反而在门徒的手前让开了。\par

\BibleQuote{我们的手也摸过生命之圣言；}即，那降于肉身者变得可触摸。正如，\par

\BibleRef{若一 1:2}\BibleQuote{生命已显现出来。}因为在福音中他这样说：\BibleQuote{凡受造的，在祂内有生命，这生命是人的光。}\BibleRef{若 1:3}\par

\BibleQuote{我们向你们显示那与圣父同在、且已向你们显现的永生。}\par

他以圣父的称呼表明圣子也无始地常存。\par

\BibleRef{若一 1:5}他说：\BibleQuote{因为天主是光。}\par

他不表达天主的本质，而是愿意宣告天主的威严，便将人眼中最美好卓越的事物用于神性。同样，保禄在谈及\Quote{不可接近的光}时也是如此。\BibleRef{弟前 6:16}而若望本人也在同一封书信中说：\Quote{天主是爱。}\BibleRef{若一 4:16}指出天主的卓越，即祂是良善仁慈的；并且因为祂是光，便借着爱德，按照灵魂的进步使人成义。因此，就本质而言不可言说的天主，就是光。\par

\Quote{在祂内毫无黑暗}——即是说，毫无情欲，对任何人都没有存留邪恶，祂不毁灭任何人，却赐给众人救恩。此外，光或指法律的诫命，或指信德，或指道理。黑暗是这些的对立面。并非因为有另一条路；因为按照天主的诫命，只有一条路。因为天主的工程是合一。二元性及所有除了合一以外的存在，都源于生活的悖逆。\par

第 7 节。他说：\Quote{祂圣子耶稣基督的血洗净我们。}\BibleRef{若一 1:7}因为主的道理，极其有力，被称为祂的血。\par

第 10 节。\Quote{如果我们说我们没有犯过罪，我们便是以祂为说谎者，祂的话就不在我们内。}\BibleRef{若一 1:10}祂的道理，或者说祂的话，就是真理。\par

第二章第 1 节。他说：\Quote{若有人犯了罪，我们在父那里有一位护慰者，耶稣基督。}\BibleRef{若一 2:1}因为主就是这样为我们作护慰者。同样，也有一位护慰者，在祂升天后，祂赏赐派遣祂。因为这些原始且首造的美德在本质上是不变的，并与从属的天使和总领天使一同施行神圣的作为——他们共享这些名称。梅瑟也以近旁和最低阶的天使之名，称呼总领天使弥额尔的美德。我们在圣先知那里也发现类似情况；但向梅瑟显现的天使是近在眼前的。梅瑟听他说，并与他面对面地交谈。对于其他先知，借着天使的中介，他们好像听到和看到一样产生了印象。\par

因此，也只有他们听见，也只有他们看见；在撒慕尔的事上也是如此。厄里叟也独自听见了那呼召他的声音。\EditorialNote{列王纪上第19章}倘若那声音是公开且普遍的，所有人都会听见。而在这一事例中，只有他一人听见，因为他内在的、由天使所产生的印象起了作用。\par

第 2 节。他说：\Quote{祂不仅为我们的罪——即信友的罪——也是为全世界的罪，}\BibleRef{若一 2:2}主是赎罪者。祂确实拯救众人；但有些人，祂借着惩罚使他们皈依而拯救他们；另一些人却自愿跟随，祂以尊荣的尊严拯救他们，以致\Quote{一切在天上的、地上的和地下的，无不向耶稣屈膝}\BibleRef{斐 2:10}，即天使、人，以及祂来临前已离开此世生活的灵魂。\par

第 3 节。\Quote{我们若遵守祂的诫命，就晓得我们认识祂。}\BibleRef{若一 2:3}因为那有知识的人（认识者）也实践属于德性领域的工作。但实践工作的人不必然是有知识的人。因为一个人可能做了正确的工作，却未必是科学奥秘的知识者。最后，知道有些工作是出于畏惧惩罚，有些是为了奖赏的许诺，他指出了因爱德而完成其工作的、被赋予知识者的完美。此外，他接着说道：\par

第 5 节。\Quote{凡遵守祂的话的，天主的爱在这人内才得以圆满：由此我们知道我们是在祂内}\BibleRef{若一 2:5}——因信德和爱德。\par

第 7 节。\Quote{我写给你们的，不是一条新诫命，而是一条旧诫命，就是你们从起初就有的}\BibleRef{若一 2:7}——即通过法律和先知；那里说：天主是唯一的。因此，他也推断说：\Quote{因为那旧诫命就是你们所听见的话。}\par

然而他又说：\par

第 8 节。\Quote{这就是那诫命；因为那黑暗}——属于悖逆，即——\Quote{已经过去，并且，看，真光已经照耀了}\BibleRef{若一 2:8}——即借着信德、借着知识、借着在人间运行的盟约、借着预备好的审判。\par

第 9 节。他说：\Quote{那说自己是在光中}——在光中，他意指在真理中——\Quote{却恨}，他说，\Quote{自己的弟兄。}\BibleRef{若一 2:9}\par

第 10 节。\Quote{那爱自己弟兄的，就住在光中，在他身上没有绊脚石。}\BibleRef{若一 2:10}\par

第 12-14 节。他随后指出仍居于肉身中的灵魂进步与成长的阶段；他把那些因主的名而被宽恕了罪的称为\Quote{孩子们}，因为许多人仅仅因着那名而相信。他把成全者称为\Quote{父老们}，\Quote{他们认识那从起初就有的}，并以理智领受了——即圣子，就是他在上面所说的\Quote{那从起初就有的。}\BibleRef{若一 2:12}\par

他说：\Quote{青年人，我写信给你们，因为你们战胜了那邪恶者。}强壮的青年人因轻视欢乐而刚强。\Quote{那邪恶者}指出魔鬼的尊高。\Quote{孩子们}则认识父；他们远离偶像，聚集于唯一的天主。\par

第 15 节。\Quote{因为世界}，他说，\Quote{是在那邪恶者的手下。}\BibleRef{若一 2:15}难道世界和世间万物不被称为天主的创造且非常好吗？是的。但是，\par

第 16 节。\Quote{肉身的贪欲、眼目的贪欲、以及人生的骄矜}——这些源于生活的悖逆——\Quote{不是出于父，而是出于世界}，也出于你们。\BibleRef{若一 2:16}\par

第 17 节。\Quote{因此，世界也要过去，连同它的贪欲；但那承行天主旨意}并遵守祂的诫命，\Quote{却永远存留。}\BibleRef{若一 2:17}\par

\BibleRef{若一 2:19} 19节。\BibleQuote{他们从我们中间出去了，但不是属于我们的}——既非背命的天使，也非堕落的人——\BibleQuote{但为要显明他们都不是属于我们的}。他清楚地区分了被选者和丧亡者的类别，以及那在信德中存留的人，\BibleQuote{从圣者领受了傅油}，这傅油是藉着信德而来的。那不在信德中存留的人。\par

\BibleRef{若一 2:22} 22节。\BibleQuote{说谎者}和\BibleQuote{假基督，就是那否认耶稣为基督的}。因为耶稣是救主和救赎者，也是基督君王。\par

\BibleRef{若一 2:23} 23节。\BibleQuote{那否认子的}，因忽视祂，\BibleQuote{就没有父，也不认识祂}。但那认识子和父的，是按着知识认识；当主在祂第二次来临时显现时，他将有自信而不会蒙羞。这种羞惭是重大的惩罚。\par

\BibleRef{若一 2:29} 29节。他说道：\BibleQuote{凡行正义的}，\BibleQuote{都是由天主而生的}；即因信德而重生。\par

\BibleRef{若一 3:1} 第三章1节。\BibleQuote{因为世界不认识我们，如同不认识祂一样}。他所说的世界，是指那些在享乐中度世俗生活的人。\par

\BibleRef{若一 3:2} 2节。\BibleQuote{亲爱的}，他说，\BibleQuote{现在我们是天主的子女}，并非出于自然亲情，而是因为天主是我们的父。因为这是更大的爱德：我们与天主本无关系，祂却依然爱我们，称我们为祂的子女。\BibleQuote{将来如何，还没有显明}；就是说，我们将达到何种荣耀。\BibleQuote{因为若是祂显现出来}——即若我们得以成全——\BibleQuote{我们将相似祂}，仿佛安息并被成义，德行纯洁，\BibleQuote{好使我们看见祂}（祂的面容）\BibleQuote{实在怎样}，是藉着领悟。\par

\BibleRef{若一 3:8} 8节。\BibleQuote{那行不义的，是属于魔鬼}，就是说，以魔鬼为父，跟随并选择同样的事。他说：\BibleQuote{魔鬼从起初就犯罪}。从那开始犯罪的起初，他不可救药地坚持犯罪。\par

\BibleRef{若一 3:9} 9节。他说：\BibleQuote{凡由天主生的，就不犯罪，因为天主的种子存留在他内}；即天主的圣言在他内，就是那藉信德重生的人。\par

\BibleRef{若一 3:10} 10节。\BibleQuote{这样，我们就知道谁是天主的子女，同样，谁是魔鬼的子女}，就是那些选择像魔鬼的事的人；因为如此他们也说属于那恶者。\par

\BibleRef{若一 3:15} 15节。\BibleQuote{凡恨自己弟兄的，就是杀人的}。因为在他内，由于不信，基督死了。因此，他接着正确地说：\BibleQuote{你们知道，凡杀人的和不信的人，没有永生存留在他内}。因为生活的基督住在信德的灵魂内。\par

\BibleRef{若一 3:16} 16节。\BibleQuote{因为祂自己为我们舍命}；就是为那些相信的人；也就是为宗徒们。如果祂为宗徒们舍命，祂指的是祂自己的宗徒们：他说的我们，就是指我们宗徒，祂为我们舍命，\BibleQuote{我们也应当为弟兄舍命}；因为邻人的救恩就是宗徒们的荣耀。\par

\BibleRef{若一 3:20} 20节。他说：\BibleQuote{因为天主比我们的心大}；就是说，天主的德能比良心大，良心将随从灵魂。因此他接着说：\BibleQuote{而且知道一切}。\par

\BibleRef{若一 3:21} 21节。\BibleQuote{亲爱的，若我们的心不责备我们，我们在天主前就有自信}。\par

\BibleRef{若一 3:24} 24节。\BibleQuote{从此我们就知道祂藉着祂赐给我们的圣神住在我们内}；就是说，藉着对将来事件的监督和预见。\par

\BibleRef{若一 4:18} 第四章18节。他说：\BibleQuote{完美的爱德驱逐恐惧}。因为信德之人的成全就是爱德。\par

\BibleRef{若一 5:6} 第五章6节。他说：\BibleQuote{这一位就是藉着水及血而来的}；并且再次——\par

\BibleRef{若一 5:8} 8节。\BibleQuote{因为作见证的有三位：圣神}，即生命，\BibleQuote{水}，即重生和信德，\BibleQuote{血}，即知识；\BibleQuote{而这三位是一致的}。因为在救主内，这些救恩的德能，生命本身存在于祂自己的圣子内。\par

\BibleRef{若一 5:14} 14节。\BibleQuote{这就是我们对祂所有的自信：如果我们按祂的旨意求什么，祂必俯听我们}。他并没有绝对地说我们将求什么，而是说我们应该求什么。\par

\BibleRef{若一 5:19} 19节。\BibleQuote{整个世界都卧在那恶者手下}；并不是指受造界，而是指世俗的人和那些按自己私欲生活的人。\par

\BibleRef{若一 5:20} 20节。\BibleQuote{天主子已经来了，赐给了我们理智}，这理智藉着信德临到我们，也称作圣神。\par

% structure: p0085 source=h2 target=section reason=relative-heading-rank; base=h2

\section{四．——关于 \BibleBook{若望}{二书} 的注释。}

\BibleBook{若望}{二书} 写给童贞女，十分简单。它是写给一位名叫厄勒克塔的巴比伦贵妇，并指示圣教会的被选。他在这书信中确立：追随信德不能没有爱德，这样就没有人能分裂耶稣基督；而只应相信耶稣基督曾降生成人。因为凡以理智领悟而拥有圣子的人，也认识圣父，并以悟性明智地把握祂那无时间起首而运作的权能的伟大。\par

\BibleRef{若二 1:10} 10节。他说：\BibleQuote{若有人来到你们那里，不带这道理，不要接他进你们家，也不要向他问安；因为向他问安的，就是有分于他的恶行}。他禁止我们向这样的人请安，并接待他们。因为对于这种人，这并非严厉。但他告诫他们，不要与那些不能明智处理神圣事物的人交谈或争辩，免得他们被似是而非的理由所影响，从而被诱离真理的道理。现在，我认为我们甚至不应与这样的人一同祈祷，因为在家中所作的祈祷中，祈祷后起身时的喜乐问安也是和平的记号。\par

% structure: p0088 source=h2 target=section reason=relative-heading-rank; base=h2

\section{二．——赫拉克肋亚主教尼基塔。}

\Emph{出自他的连珠集}\par

\Emph{一、——\BibleRef{约 1:21}}\par

但约伯的话可以更优雅地理解为关于邪恶和罪，即：\Quote{赤身}最初从大地形成，仿佛出自\Quote{母胎}；\Quote{我也要赤身归去}；赤身，不是指财物，因为那是琐碎而平常的事，而是指邪恶和罪，以及那些生活败坏的人所追随的丑陋形态。显然，我们所有人都是赤身而生，又赤身被埋葬，只裹着寿衣。因为天主为我们预备了另一种生命，并使现世生命成为通往那生命的道路；祂将我们从所拥有的财货中获得的供给指定为路上的干粮；而当我们离开这条路时，由我们所拥有的东西构成的财富就不再与我们同行。因为我们所拥有的，没有一样真正属于我们自己：唯独一种拥有，即虔敬，才是我们真正的所有物。死亡临到时，不会夺去这虔敬；但死亡会将我们从其他一切中逐出，尽管我们不愿意。因为我们所拥有的，都是为了维持生命而领受的；每个人仅仅享受了它们的使用之后便离去，从生命中只得到短暂的记忆。因为这就是一切繁荣的终结；这就是现世美好事物的结局。因此，婴儿在离开母胎、睁开眼睛时，立即开始哭泣，而不是欢笑。因为他哭泣，仿佛哀悼生命，生命从一开始就让他尝到死亡之礼。因为一出生，他的手和脚就被包裹起来；被束缚着吃奶。啊，生命的开端，死亡的前兆！孩子刚进入生命，就立刻穿上了死者的衣袍：因为大自然提醒出生的人他们的结局。因此，孩子出生时哭喊，仿佛向母亲哀诉。母亲啊，你为什么把我带到这个生命里来，在这里生命的延长就是走向死亡？你为什么把我带到这个动荡的世界，在这里一出生，襁褓就是我最初的经历？你为什么把我交付给这样一个生命，在其中，可悲的青春在衰老之前就消逝，而老年又被躲避，如同在死亡的判决之下？母亲啊，生命的历程是可怕的，它以死亡为奔跑者的目标。我们行走的生命之路是苦涩的，坟墓是行旅的客栈。我们航行的生命之海是危险的，因为冥府如同海盗攻击我们。唯独人生来完全赤身，没有与生俱来的武器或衣服；并非低于其他动物，而是要使赤身和一无所携激发思考；思考引出灵巧，驱除懒惰，引入满足需求的技艺，并产生各种发明。因为赤身的人充满巧计，被需要如同被刺棒驱策：如何躲避雨水，如何逃避寒冷，如何抵挡打击，如何耕种土地，如何恐吓野兽，如何制服那些更强大的野兽。被雨淋湿，他构思了屋顶；受冷后，他发明了衣服；被击打后，他制造了胸甲；耕种时荆棘刺破双手，他借助工具；赤身易受野兽捕食，他从恐惧中发现了恐吓恐惧者的技艺。赤身产生了一个又一个成就；因此，他的赤身也成了一项恩赐和主人般的恩惠。因此，约伯在失去财富、财产、子女的祝福、众多的后裔，并在短时间内失去一切之后，发出了这感恩的呼喊：\Quote{我赤身出于母胎，也必赤身归回那里}——即归回天主，归回那有福的份和安息。\par

% structure: p0092 source=h2 target=section reason=relative-heading-rank; base=h2

\section{二．——出自同一著作}

\BibleRef{约 5:7}。宁静是所有事物中最值得珍视的。为此，圣言向我们提出了有福的约伯作为榜样。因为论到他经上记着说：\Quote{有谁像约伯，像喝水一样吞下嘲笑呢？}因为一个人若能达到如此程度的含忍，以至于能轻松地应对因受冤屈而生的真正尖锐且不易被所有人征服的痛苦，那真是令人羡慕，在我看来，值得一切钦佩。\par

% structure: p0094 source=h2 target=section reason=relative-heading-rank; base=h2

\section{三．——出自尼塞塔斯《玛窦福音连珠集》}

\BibleRef{玛 5:42}。施舍应当有判断，给与配得的人，好让我们从至高者那里获得赏报。但那些拥有却以虚假借口索取的人，或者那些能够自助却想从他人那里获取的人是有祸的。因为那拥有却为了行骗或出于懒惰而索取的人，将被定罪。\par

% structure: p0096 source=h2 target=section reason=relative-heading-rank; base=h2

\section{四．——出自同一著作}

\BibleRef{玛 13:31}。宣告天国的圣言就像芥子一样尖锐辛辣，能抑制胆汁（即愤怒），能平息炎症（即骄傲）；从这圣言中涌出灵魂真正的健康和永恒的健全。圣言的生长竟达到如此大的规模，以至于从它生出的树（即遍及全地的基督的教会）充满了世界，以致天空的飞鸟——即神圣的天使和高尚的灵魂——栖居在它的枝上。\par

% structure: p0098 source=h2 target=section reason=relative-heading-rank; base=h2

\section{五．——出自同一著作}

\BibleRef{玛 13:46}。一颗珍珠，透明而光芒纯净，就是耶稣，童贞女所生的神性闪电。因为正如珍珠在肉体、蚌壳和水分中产生，显现为湿润透明、充满光明和灵气的物体；同样，天主圣言降生成人，是理智的光，透过光明而湿润的身体放射光芒。\par

% structure: p0100 source=h2 target=section reason=relative-heading-rank; base=h2

\section{三．——出自科尔德里乌斯所编《路加福音连珠集》}

\BibleRef{路 3:22}。天主在此取了\Quote{形状}，不是人的形状，而是\Quote{鸽子的形状}，因为祂愿意藉着圣神以鸽子的形状显现的新方式，宣告祂的纯朴和威严。\par

\BibleRef{路 16:17}。或许藉着一划和一钩，祂的公义呼喊说：\Quote{你若正直地到我这里来，我也必正直地到你那里去；但你若弯曲，我也必弯曲地来，万军的上主说。}暗示罪人的道路是迂回弯曲的。因为由耶稣的一划所暗示的正直而合乎本性的道路，就是祂的良善，这良善经常引导那些因听见而相信的人。\Quote{所以，法律上的一一划或一钩也不能过去。}即无论是关于正直和良善的相互应许，还是关于弯曲和不义者所应得的惩罚，都不能过去。\Quote{因为上主善待善人，但那些转向弯曲道路的人，上主要引导他们与作恶的人同行。}\par

% structure: p0103 source=h2 target=section reason=relative-heading-rank; base=h2

\section{IV.— 出自 \WorkTitle{要义} 诸书。}

% structure: p0104 source=h2 target=section reason=relative-heading-rank; base=h2

\section{Œcumenius 从第三卷，论格前十一10。}

\Quote{为了天使的缘故。}他所说的天使是指义人和有德的人。那么，她应当蒙着头，免得使她绊倒他们而陷于淫乱。因为天上的真天使即使在她蒙着头时也看得见她。\par

% structure: p0106 source=h2 target=section reason=relative-heading-rank; base=h2

\section{同卷，第四卷，论格后五16。}

\Quote{即使我们曾按肉性认识过基督。}正如在我们身上，\Quote{按肉性}是指处于罪恶之中，而脱离罪恶便是\Quote{不按肉性}；同样，在基督身上，\Quote{按肉性}是指祂受制于自然情感，而不受制于这些情感便是\Quote{不按肉性}。但他说，正如祂被释放了，同样我们也被释放了。\par

% structure: p0108 source=h2 target=section reason=relative-heading-rank; base=h2

\section{同卷，第四卷，论格后六11。}

\Quote{我们的心敞开了，}为了教导你们一切。但你们在自己的心肠中受到限制，也就是说，在对天主的爱中受到了限制，而你们本应在这爱中爱我。\par

% structure: p0110 source=h2 target=section reason=relative-heading-rank; base=h2

\section{同卷，第五卷，论迦五24。}

\Quote{凡属于基督的人，已把肉身钉在十字架上了。}何必再一一列举德行的各个方面呢？因为有些人已将自己钉在十字架上，就情欲而言，也把情欲钉在了自己身上。按照这一解释，\Quote{和}字并非多余。\Quote{凡属于基督的人}——即追求祂的人——\Quote{已把自己的肉身钉在十字架上。}\par

% structure: p0112 source=h2 target=section reason=relative-heading-rank; base=h2

\section{Moschus：\WorkTitle{灵性草地}，第五卷，第176章。}

是的，确实，宗徒们受了洗，正如克莱孟在 \WorkTitle{要义} 第五卷所述。因为他在解释宗徒的话\Quote{我感谢天主，我没有给你们任何人付洗}时说，据说基督只给伯多禄付了洗，伯多禄给安德肋付了洗，安德肋给若望付了洗，而他们给雅各伯和其余的人付了洗。\par

% structure: p0114 source=h2 target=section reason=relative-heading-rank; base=h2

\section{Eusebius：\WorkTitle{教会史}，第六卷，第二章第一节。}

如今克莱孟在 \WorkTitle{要义} 第六卷中作了如下陈述。他说，伯多禄、雅各伯和若望在救主升天后，虽蒙主特别尊荣，却不争夺荣耀，反而立了义人雅各伯为耶路撒冷的主教。\par

% structure: p0116 source=h2 target=section reason=relative-heading-rank; base=h2

\section{Eusebius：\WorkTitle{教会史}，第二卷，第十五章。}

于是，藉着天主圣言对他们的造访，西满的能力被消灭了，并立即与他本人一同毁灭。如此虔诚的光辉照耀在伯多禄听众的心上，使他们不以一次听闻或天神圣传的口头教导为满足，而是用各种恳求坚请那位福音所归功的人——即伯多禄的同伴马尔谷——将口头传授给他们的教导用文字记录下来；他们一直缠着这个人直到说服了他；因此，我们亏欠他们那部称为\BibleBook{马尔谷}{福音}的圣经。据说，宗徒藉着圣神的启示得知所发生的事后，对人们的热情感到高兴，并批准了这部著作供教会诵读。克莱孟在 \WorkTitle{要义} 第六卷中叙述了此事。\par

% structure: p0118 source=h2 target=section reason=relative-heading-rank; base=h2

\section{Eusebius：同上。}

然后，正如圣经所载，黑落德在处死雅各伯后，见所作的事使犹太人欢喜，便也下手捉了伯多禄；将他锁在狱中，本要立刻处死他，若不是一位天使在夜间以异象向他显现，奇妙地解开锁链，派他去从事宣讲的职务。\par

% structure: p0120 source=h2 target=section reason=relative-heading-rank; base=h2

\section{Eusebius：\WorkTitle{教会史}，第六卷，第十四章。}

并且在 \WorkTitle{要义} 中，他简要地概括了整个约书的圣经；并且没有跳过有争议的书卷——我指的是 \BibleBook{犹达}{书} 和其余的公函，以及 \WorkTitle{巴尔纳伯书}，和所谓的 \WorkTitle{伯多禄启示录}。他还说，\BibleBook{希伯来}{书} 是保禄所写，并且是用希伯来语写给希伯来人的；但路加仔细地翻译了它，把它给了希腊人，因此在这封书信和 \BibleBook{}{宗徒大事录} 中可以发现相同的表达色彩；并且 \Quote{保禄宗徒} 这个名衔很适当地没有放在开头，因为，他说，写信给对保禄有偏见并怀疑的希伯来人，他极其智慧地没有一开始就写上自己的名字来排斥他们。\par

% structure: p0122 source=h2 target=section reason=relative-heading-rank; base=h2

\section{Eusebius：第七卷。}

\BibleRef{弟前 2:6}。\Quote{在适当的时候；}即当人们处于适合信仰的状态时。\par

\BibleRef{弟前 3:16}。\Quote{被天使看见。}啊，奥妙！天使在基督与我们同在时看见了祂，而在这之前他们未曾见过祂。不像人那样。\par

\BibleRef{弟前 5:8}。\Quote{尤其是对自己家中的人。}那照顾自己家中的人，不仅照顾亲属，而且也藉着根除情欲照顾自己的人，才是照顾自己的家人。\par

\BibleRef{弟前 5:10}。\Quote{如果她洗过圣徒的脚；}即如果她曾毫无羞耻地为圣徒做过最卑微的服事。\par

\BibleRef{弟前 5:21}。\Quote{不存偏见；}即不因任何错误步骤而陷入不顺服的审判和惩罚之下。\par

\BibleRef{弟前 6:13}。\Quote{在般雀比拉多前作过见证。}因为祂藉着祂所做的事见证祂是基督、天主子。\par

\BibleRef{弟后 2:2}。\Quote{由许多见证人；}即法律和先知。因为这些是宗徒为自己宣讲所作的见证人。\par

% structure: p0130 source=h2 target=section reason=relative-heading-rank; base=h2

\section{Eusebius：\WorkTitle{教会史}，第七卷，第二章第一节。}

主在复活后，将\OriginalTerm{知识}{\GreekText{τὴν} \GreekText{γνῶσιν}}传授给了义人雅各伯、若望和伯多禄。他们将此传授给了其余宗徒，其余宗徒传授给了七十人，巴尔纳伯是其中之一。\par

% structure: p0132 source=h2 target=section reason=relative-heading-rank; base=h2

\section{Eusebius：同上，第二章第二节。}

有关这位雅各伯，克莱孟在\WorkTitle{纲要}第七卷中也记载了一件值得纪念的轶事，来源自其前辈的传统。他说，那带他受审的人，看见他为主作证，深受感动，便承认自己也是基督徒。于是，两人一同被带走，路上那人请求雅各伯原谅他。雅各伯沉思片刻，说：\Quote{愿你平安}，并亲吻了他。就这样，两人一同被斩首。\par

% structure: p0134 source=h2 target=section reason=relative-heading-rank; base=h2

\section{欧瑟伯：同上，第六卷第十四章}

有福的司铎曾说过：既然主作为全能者的宗徒被派遣到希伯来人那里，那么保禄作为被派遣到外邦人的人，出于谦虚和对主的敬畏，没有自称为希伯来人的宗徒；因为他既是外邦人的宣讲者和宗徒，那么写信给希伯来人便超出了他本来的职责。\par

% structure: p0136 source=h2 target=section reason=relative-heading-rank; base=h2

\section{欧瑟伯：同上}

此外，在同一本书中，克莱孟记载了他从前辈长老那里接受的传统，关于福音书的次序，内容如下。他说，包含族谱的福音书写得最早，而\WorkTitle{马尔谷福音}是在以下情形中写成的：\par

伯多禄在罗马公开宣讲圣言，并借着圣神宣报了福音；在场的人很多，他们恳请马尔谷——因为他从早期就一直跟随伯多禄，并且记住了所说的话——将所传讲的内容记录下来。他写成福音后，交给了那些请求他的人；此事传到伯多禄耳中，他既不阻止也不鼓励。最后，若望，看见福音书中已记载了有形的事物，便在密友的请求下，受圣神感动，写成了一部属灵的福音。\par

% structure: p0139 source=h2 target=section reason=relative-heading-rank; base=h2

\section{五、摘自\WorkTitle{论天主眷顾}}

% structure: p0140 source=h2 target=section reason=relative-heading-rank; base=h2

\section{圣玛西穆，第二卷第一百一十四页}

存有在天主内。天主是神圣的存有，永恒无始，无形无体，无限无量，是万有之因。存有就是完全自立的。自然是事物的真理，或事物的内在实在。另一些人认为，自然是已经存在之物的产生。又有人认为，自然是天主的眷顾，它导致受造物的存有和存有的方式。\par

% structure: p0142 source=h2 target=section reason=relative-heading-rank; base=h2

\section{圣玛西穆：同上，第一百五十二页}

意愿是一种自然的能力，它渴望符合自然的事物。意愿是理性的受造物与本性相应的自然欲望。意愿是自决心灵的一种自然自发的运动，或者是心灵自愿地朝向任何事物的运动。自发性是自然运动的心灵，或是灵魂的理智性的自决运动。\par

% structure: p0144 source=h2 target=section reason=relative-heading-rank; base=h2

\section{六、摘自\WorkTitle{论灵魂}}

% structure: p0145 source=h2 target=section reason=relative-heading-rank; base=h2

\section{玛西穆与安多尼·梅利撒}

那些从一切事物中自由呼吸的灵魂拥有生命，虽然与身体分离，却依然怀着对身体的渴望，不朽地被带到天主的怀抱中；犹如冬季里，大地的蒸汽被阳光吸引而上升到太阳那里。\par

% structure: p0147 source=h2 target=section reason=relative-heading-rank; base=h2

\section{巴罗克手抄本}

所有灵魂都是不朽的，甚至恶人的灵魂也是如此；对他们来说，不死反而更糟糕。因为他们被不灭之火的永恒惩罚所折磨，却又不会死去，所以他们的痛苦不可能有终结。\par

% structure: p0149 source=h2 target=section reason=relative-heading-rank; base=h2

\section{七、摘自\WorkTitle{论毁谤}}

% structure: p0150 source=h2 target=section reason=relative-heading-rank; base=h2

\section{安多尼·梅利撒，第二卷，第六十九讲}

绝不要害怕那向你说话的毁谤者。反而要说：\Quote{住口吧，兄弟；我每天犯下更严重的过错，又怎能判断他呢？}因为你将获得两样益处：用同一帖膏药治愈你自己和你的邻人。毁谤者显明了什么是真正的邪恶。因此，天主通过这些辩论，使每个人的性情得以显明。\par

% structure: p0152 source=h2 target=section reason=relative-heading-rank; base=h2

\section{安多尼·梅利撒，第一卷，第六十四讲，及第二卷，第八十七讲。亦见玛西穆，第五十九讲，第六百六十九页；大马士革的若望，第二卷。}

使信徒成义的，不是禁戒行为，而是思想的纯洁和真诚。\par

% structure: p0154 source=h2 target=section reason=relative-heading-rank; base=h2

\section{八、摘自安多尼·梅利撒的其他残篇}

% structure: p0155 source=h2 target=section reason=relative-heading-rank; base=h2

\section{一、第一卷，第十七讲，论告解}

当灵魂的过失一发生，悔改便立即接纳，不拖延，不让冲动持续长时间，那么悔改就能擦除一切罪过。因为这样，邪恶就无法在我们身上留下痕迹，它在攻击的那一刻就被拔除，犹如一棵新栽的植物。\par

正如被称为螃蟹的生物，由于时而前进时而后退而容易被捕捉；同样，灵魂时而欢笑，时而哭泣，时而耽于逸乐，便无法行善。\par

那有时悲伤、有时快乐发笑的人，就像一个人用面包投掷 \OriginalTerm{纵欲}{voluptuousness} 的狗，表面上是在赶它，实际上是在邀请它留在自己身边。\par

% structure: p0159 source=h2 target=section reason=relative-heading-rank; base=h2

\section{二、第一卷，第五十一讲，论赞美}

一些谄媚者向一位智者道贺。他对他们说：\Quote{如果你们停止赞美我，你们走后我还以为自己很伟大；但如果你们不停止赞美我，我就猜到自己有多不洁。}\par

虚假的赞美不如真实的责备有价值。\par

% structure: p0162 source=h2 target=section reason=relative-heading-rank; base=h2

\section{三、第二卷，第四十六讲，论懒惰和怠惰者}

对软弱和病弱的人而言，适度也显得过度。\par

% structure: p0164 source=h2 target=section reason=relative-heading-rank; base=h2

\section{四、第二卷，第五十五讲，论你的近人——当你应背负他的重担等}

出于知识的责备是十分可靠的。有时，被定罪之人的知识也被证明是最完美的证明。\par

% structure: p0166 source=h2 target=section reason=relative-heading-rank; base=h2

\section{五、第二卷，第七十四讲，论骄傲及贪图虚荣者}

那自高自大的人，迅速堕落到卑微的地位，正如圣言所教导的。\par

% structure: p0168 source=h2 target=section reason=relative-heading-rank; base=h2

\section{六、第二卷，第八十七讲}

纯洁的言语和无瑕的生活是天主的宝座和真正的圣殿。\par

% structure: p0170 source=h2 target=section reason=relative-heading-rank; base=h2

\section{九、\WorkTitle{论婚姻}的残篇}

% structure: p0171 source=h2 target=section reason=relative-heading-rank; base=h2

\section{玛西穆，第三讲，第五百三十八页，论谦逊与贞洁。亦见大马士革的若望，第三卷——平行第二十七章。}

不独淫乱被称为淫乱，过早结婚也被称为淫乱；也就是说，当一个人尚未成熟就被嫁给丈夫，无论出于自愿还是父母的安排。\par

% structure: p0173 source=h2 target=section reason=relative-heading-rank; base=h2

\section{十、其他失传著作的残篇}

% structure: p0174 source=h2 target=section reason=relative-heading-rank; base=h2

\section{玛西穆，第二讲。——大马士革的若望，第二卷第七十章。——安多尼·梅利撒，第一卷第五十二讲。}

谄媚是友谊的毒药。大多数人习惯于讨好君王的运气，而不是君王本人。\par

% structure: p0176 source=h2 target=section reason=relative-heading-rank; base=h2

\section{玛克西谟，\WorkTitle{讲道集}第十三篇，第574页。— 安多尼·梅利莎，\WorkTitle{讲道集}第三十二篇，第45页，及\WorkTitle{讲道集}第三十三篇，第57页。}

喜爱节俭的人躲避奢侈，视之为灵魂和身体的祸害。拥有并使用必需品，在性质上并无害处，但在数量上若超过分寸则有害。食物匮乏是一种必要的益处。\par

% structure: p0178 source=h2 target=section reason=relative-heading-rank; base=h2

\section{玛克西谟，\WorkTitle{讲道集}第五十二篇，第654页。— 安多尼·梅利莎，\WorkTitle{讲道集}第一卷，第五十四篇。}

对死亡的真切记忆是饮食的节制；当饮食减少时，情欲也随之减少。\par

% structure: p0180 source=h2 target=section reason=relative-heading-rank; base=h2

\section{玛克西谟，\WorkTitle{讲道集}第五十五篇，第661页。}

最重要的是，基督徒不允许以暴力纠正罪恶过犯。因为天主加冕的不是那些因强迫而避免邪恶的人，而是那些因选择而避免邪恶的人。一个人不可能持续为善，除非出于他自己的选择。因为因他人强迫而成为善的人并非善人；因为他不是出于自己的选择而成为善的。正是每个人的自由意志造就了真正的善，也揭示了真正的恶。因此，通过这些倾向，天主设法使祂自己的倾向得以显明。\par

% structure: p0182 source=h2 target=section reason=relative-heading-rank; base=h2

\section{十一、仅见于牛津版希腊文残篇}

% structure: p0183 source=h2 target=section reason=relative-heading-rank; base=h2

\section{出自\WorkTitle{最后关于逾越节的著作}}

因此，在过去的年代里，耶稣曾去享用犹太人所祭献的逾越节晚餐，遵守庆节。但是，当祂宣讲了那原是逾越节的羔羊、被牵到宰杀之地如绵羊的天主羔羊后，随即在第十三日教导门徒这预像的奥秘，正是在这一天他们询问说：\Quote{祢要我们在哪里为祢预备吃逾越节晚餐？}\BibleRef{玛 26:17}所以，正是在这一天，祝圣无酵饼和预备庆节的事都完成了。因此，\Person{若望}自然描绘门徒早已预备好，让主为他们洗脚。次日，我们的救主——那被犹太人作为赎罪祭献的逾越节羔羊——受苦受难。\par

% structure: p0185 source=h2 target=section reason=relative-heading-rank; base=h2

\section{同一著作}

因此，与第十四日相宜，在祂受苦的当天清晨，司祭长和经师们将祂带到比拉多那里，却没有进入\OriginalTerm{总督府}{Prætorium}，以免沾染不洁，好能自由地在晚间吃逾越节晚餐。对于这一精确的日子，全部圣经都一致赞同，福音书也彼此和谐。复活也为此作证。祂确在第三日复活了，这日正是收获周的第一日，法律规定司祭应在该日奉献禾捆。\par

% structure: p0187 source=h2 target=section reason=relative-heading-rank; base=h2

\section{玛加略·金首：\WorkTitle{浪子回头比喻}，\BibleBook{路加}{福音}第十五章，\WorkTitle{论路加福音第十五章}讲道，接近结尾}

1. 在天上，若有人成为从圣父管治下生活的流亡者和逃遁者，不知道那些远离祂的人必要灭亡；若有人挥霍了圣父的恩赐、家产和产业；若有人信德败坏、望德耗尽，因与外邦人一同奔入同样的放荡淫逸；然后，饥饿穷乏，甚至不能以猪所吃的充饥，却起来回到他父亲那里，那是何等歌舞和高举的庆节啊！\par

然而，慈祥的父亲并不等儿子来到他跟前。因为若非他感到父亲的仁慈，他或许永远不能也不敢靠近。因此，当他刚有愿望，当他立即开始行动，当他迈出第一步，还离得远的时候，祂（圣父）就动了怜悯的心，跑去扑在他颈上，亲吻了他。于是儿子鼓起勇气，承认了自己的所作所为。\par

因此，圣父将当得的荣耀和尊荣赐给他，给他穿上上好的袍子，即不朽的袍子；并戴上戒指，即君王的印玺和神圣的印章——祝圣印记、荣耀签名、见证抵押（因为经上说：\Quote{天主是真实的，那人已在他身上盖了印}\BibleRef{若 3:33}），还有鞋子，不是那双会朽坏的、他踏在圣地上时被命脱下的鞋子，也不是那双被禁止穿去宣讲天国福音的鞋子，而是穿不坏、适合天国旅程的鞋子，适宜并装饰天路，未洗过的脚从不穿上，只有经我们师主洗过的脚才能穿。\par

罪人的灵魂有许多鞋子，被它们捆绑和束缚。因为各人被自己的罪恶绳索所捆绑。因此，\Person{亚巴郎}向\Place{索多玛}王起誓说：\Quote{凡属于你的，连一根线、一根鞋带，我都不拿。}\BibleRef{创 14:23}正因为这些鞋子在地上被玷污和污染，各种不义和自私占据了人生。正如主通过\Person{亚毛斯}责备以色列说：\Quote{为了以色列的三重罪过，甚至四重罪过，我决不收回成命，因为他们为了银钱出卖义人，为一双鞋出卖穷人，用脚践踏地上的尘土。}\BibleRef{亚 2:6}\par

2. 如今，圣父吩咐仆人给那悔改并投奔祂的儿子拿去的鞋，既不阻碍也不拖累大地（因为地上的帐幕压住焦虑的心思）；而是轻盈、上升、飘向天堂，并为那已转向圣父的人充当所需的梯子和战车。因为，在外部被这一切美好地装饰之后，他进入内在的喜乐。因为\Quote{拿来}是那先前说过\Quote{他离得还远时，就跑去扑在他的颈上}者所说的。因为在这里，我们必须完成进入那被邀请的婚宴的一切准备。于是，那已准备好进入的人要说：\Quote{我的这喜乐满足了。}\BibleRef{若 3:29} 但不可爱、不体面的人会听见：\Quote{朋友，你怎么没有婚宴礼服就进来了？}\BibleRef{玛 22:12} 那肥美丰腴的食物——丰盛充足的佳肴——那肥牛犊被宰了；它也被称为羔羊（并非字面意思），以免有人认为它渺小；但它伟大又最伟大。因为\Quote{天主的羔羊，除免世罪者}并非渺小，\BibleRef{若 1:29} 祂\Quote{被牵去如同待宰的羊}，那充满骨髓的祭品，其所有脂肪，按神圣的法律，皆属上主。因为祂完全献身并奉献给上主，如此健壮，如此巨大，以至于遍及一切，并使吃祂、以祂为粮的人得到饱足。因为祂既是肉，也是饼，并将自身作为两者赐给我们食用。\par

于是，对于前来投奔祂的儿子们，圣父赐下牛犊，它被宰杀并食用。但那些不来投奔祂的人，祂追赶并剥夺继承权，并显为极有力的公牛。此处，因祂的伟岸和勇力，论及祂说：\Quote{祂的荣耀如同独角兽的荣耀。}\BibleRef{户 23:22} 哈巴谷先知看见祂带着角，并赞美祂的防御姿态：\Quote{角在祂手中。}\BibleRef{哈 3:4} 因此，这记号显示祂的权能和权威——双面穿刺的角，或者更确切地说，刺透一切。那些吃的人如此得着力量，从内在的生命之粮中保持如此强大的力量，以至于他们自己比敌人更强壮，几乎装备了公牛的角，正如经上说：\Quote{因你，我们将抵撞我们的仇敌。}\par

3. 那里有喜乐、音乐和舞蹈；尽管那长久与圣父同在、常顺服圣父的长子，看到有罪之人得享幸福，便心里不快。\par

于是圣父叫他，说：\Quote{儿子，你常与我同在。}还有什么比与天主常在、侍立祂前、事奉祂更大的喜乐、盛宴和庆典呢？\Quote{凡我所有的，都是你的。}那承受天主产业的人是有福的，圣父为他保有产业——忠信之人，万物都属他。\par

\Quote{我们理当欢喜快乐，因为你的兄弟死而复活了。}仁慈的圣父，祢赐万物生命，并使死者复活。\Quote{他迷失了，又找到了。}\Quote{祢所拣选并接纳的人是有福的}，祢寻找他，便找到了他。\Quote{得赦免其过、遮盖其罪的人是有福的。}人的本分是悔改己罪，但愿这悔改伴随着不可阻遏的改变。因为不如此行的人必蒙羞，因他并非全心行事，而是草率从事。\par

我们当逃向天主。我们要不懈而奋力地追求此事。因为祂说：\Quote{凡劳苦和负重担的，你们都到我这里来，我要使你们安息。}\BibleRef{玛 11:28} 祈祷和谦卑的告解是自愿的行为。因此命令道：\Quote{先陈述你的罪，好使你成义。}\BibleRef{依 43:26} 此后我们将获得什么，我们将成为什么，并非由我们判断。\par

4. 这是比喻的严格含义。悔改的儿子来到怜悯的圣父面前，从未奢望这些——上好的袍子、戒指和鞋，或品尝肥牛犊，或分享喜乐，或享受音乐和舞蹈；他只要得到自认为配得的东西就满足了。他曾决心说：\Quote{把我当作祢一个雇工吧。}但当他看到圣父的欢迎迎接他时，他没有说这话，而是先说出心中所想：\Quote{圣父，我得罪了天，也得罪了祢。}于是，他的谦卑和自责成了成义和光荣的原因。因为义人在第一句话中便责备自己。同样，税吏离开时比法利塞人更得成义。因此，儿子既不知道他将获得什么，也不知道如何接受、使用或穿上所赐之物；因为他没有自己拿起袍子穿上。而是说：\Quote{给他穿上。}他没有自己把戒指戴在手指上，而是那些奉命的人\Quote{把戒指戴在他手上。}他也没有自己穿上鞋，而是那些听从的人\Quote{把鞋穿在他脚上。}\par

这些事对他人或对他自己而言，也许难以置信，在发生之前也出乎意料；但赠予他的礼物被欣然接受并受到赞美。\par

5. 这比喻表明：运用理性的能力已赐予每个人。因此，荡子被描述为向父亲索要他的那份家产，即理性赋予的心智状态。因为理性之拥有被赐予所有人，为追求善而避免恶。但许多由天主装备了这理性者，却滥用所赐的知识，堕入恶行的放荡，邪恶地浪费理性的实质——眼目用于可耻的景象，口舌用于亵渎的话语，嗅觉用于污秽淫逸的过度享乐，口腹用于猪般的贪食，双手用于偷窃，双脚用于奔入阴谋，心思用于不敬的计谋，倾向用于对安逸之爱的纵容，头脑用于野蛮的消遣。他们丝毫不保留理性的实质不被挥霍。因此，基督在比喻中描绘了这样一个人——一个理性的受造物，理性昏暗，向天主祈求合乎理性之事；然后从天主获得，却邪恶地使用所赐之物，尤其是已赐予他的圣洗恩惠；因此祂也称他为荡子；然后，在所赐之物被挥霍之后，又因悔改而恢复，[祂描绘了]天主向他显示的爱。\par

6. 因为祂说：\BibleQuote{把那肥牛犊牵来，宰了，让我们吃喝快乐；因为我这个儿子}——这是最近亲的称谓，并指示了赐予信徒的恩赐——\BibleQuote{是死而失丧的}，这是最疏远的表达；因为对于活着的人，还有什么比失丧和死亡更疏远呢？因为两者都不能再被拥有。但从最近亲的关系堕落到最疏远的境地，又因悔改回到近亲的关系。因为经上说：\BibleQuote{给他穿上最好的袍子}，那是他领受圣洗时即刻拥有的。我指的是圣洗的荣耀、罪的赦免，以及其他恩惠的分享，他一旦触碰圣洗池就获得了。\par

\BibleQuote{把戒指戴在他手上。}这是天主圣三的奥迹；这是印在信徒身上的印记。\par

\Quote{把鞋穿在他脚上}，作为\BibleQuote{和平福音的准备}\BibleRef{弗 6:15}，以及走向善行的整个道路。\par

7. 但基督在圣洗后犯罪的人中寻找失丧者，而天主的敌人\Person{诺瓦替安}却将他们弃于毁灭。因此，我们若悔改，就不要认为任何过犯不可挽回；警惕跌倒，若已跌倒，就回头。既畏惧犯罪，犯罪后也要避免绝望，并渴望坚定；沉沦时，要赶紧起来。让我们听从主，祂召唤我们：\BibleQuote{凡劳苦和负重担的，你们都到我跟前来，我要使你们安息。}\BibleRef{玛 11:28}让我们运用理性的恩赐于明智的行为。让我们现在学习戒避邪恶，免得将来被迫学习。让我们把生命当作善行的训练场；让我们被唤醒去憎恨罪恶。让我们怀着对造物主的深爱；让我们全心依恋祂；不要像荡子那样邪恶地浪费理性的实质。让我们获得那积存的喜乐，\Person{保禄}曾欢呼说：\BibleQuote{谁能使我们与基督的爱隔绝？}\BibleRef{罗 8:35}愿荣耀和尊崇归于祂，及于圣父和圣神，世世无穷。阿们。\par

% structure: p0205 source=h2 target=section reason=relative-heading-rank; base=h2

\section{\Person{玛加略·克律索克法洛斯}：第八篇讲道。论\BibleBook{玛窦}{福音}第八章，以及第七卷。论\BibleBook{路加}{福音}第十三章。}

因此，天主在此并非取人的形状，而是取鸽子的形状，因为祂愿意借鸽子的相似，显示圣神新显现的纯朴与温柔。因为法律是严厉的，用刀剑惩罚；但恩宠是喜乐的，用温和之言训练。因此，主也对那些说祂应当像\Person{厄里亚}那样，用火惩罚不接待祂的人的宗徒们说：\BibleQuote{你们不知道你们是什么精神。}\BibleRef{路 9:55}\par

% structure: p0207 source=h2 target=section reason=relative-heading-rank; base=h2

\section{选自同书——第十三卷，第九章。}

或许藉着\Quote{一撇和一捺}，祂的正义呼喊：\BibleQuote{若你们正直地待我，我也必正直地待你们；若你们行事乖僻，我也必行事乖僻，}这是万军的上主说的\BibleRef{肋 26:24}，意指以乖僻之道指称罪人的过犯。因为正直的、合乎本性的道路，由耶稣的一撇所指明的，是祂的良善，对于顺从相信者坚定不移。因此，法律的一撇或一捺都不会废去；既不会废去应许给正直行路者的应许，也不会废去对偏离者所威胁的惩罚。因为主善待行正直道路的人；但\BibleQuote{那些随从乖僻道路的，祂必引他们与作恶的人一同出去。}\BibleQuote{与纯全者，祂纯全；与乖僻者，祂乖僻；}对乖僻者，祂降下乖僻的道路。\par

天主将祂自己光明的肖像，即最伟大的那一位，像印记般印在按祂相似受造的人身上，使人能成为万物的统治者和主宰，并使万物服事他。因此，天主判定人完全是祂的，是祂自己的肖像。祂是不可见的；但祂的肖像——人是可见的。因此，人无论对他人行善或作恶，都归到祂身上。因此，审判将从祂而出，按各人的功过指定一切；因为祂要为自己肖像复仇。\par

% structure: p0210 source=h2 target=section reason=relative-heading-rank; base=h2

\section{十二、牛津版未收录的残篇。}

% structure: p0211 source=h2 target=section reason=relative-heading-rank; base=h2

\section{一、见于亚纳斯大西·西乃，《疑问集》第96问。}

正如人现在仍能按亚当原初的构造造人，天主不再创造，因为祂一次而永远地赐予人生育人的能力，对我们的本性说：\BibleQuote{你们要生育繁殖，充满大地。}\BibleRef{创 1:28}同样，祂也以其全能全知的大能，安排我们身体的分解与死亡按自然顺序和法则发生，通过其元素的转化，符合祂神圣的知识与理解。\par

% structure: p0213 source=h2 target=section reason=relative-heading-rank; base=h2

\section{二、\Person{君士坦丁堡宗主教若望·贝库斯}论圣神的发源。载于\Person{利奥·阿拉提乌斯}\WorkTitle{著作集}卷一，第 248 页。}

再者，斯卓玛提人克莱孟在他所制定的各种定义中，旨在引导有志于研习神学的人明白宗教的每一条教义，在定义圣神是什么、如何被称作圣神时，说：\Quote{圣神是一种实体，精微、非物质，且无形地流出。}\par

% structure: p0215 source=h2 target=section reason=relative-heading-rank; base=h2

\section{三、选自君士坦丁堡的尼塞福鲁斯未刊行的\WorkTitle{反圣像破坏者论辩}，由勒努里在其\WorkTitle{教父文库导论}中以希腊文和拉丁文出版，卷一，第 1334 a.b 页。选自亚历山大城司铎克莱孟所著\WorkTitle{反犹太化者之书}。}

达味之子撒罗满在题为 \BibleBook{}{列王纪} 的书中，不仅指明真圣殿的结构是属天属灵的，也关联到那既是达味之子又是达味之主的祂将要建造的肉身，为祂自己的临在——祂决意在此作为活像建立自己的居所——以及通过信德的合一而兴起的教会，他明确地说：\BibleQuote{天主果真与人住在地上吗？}\BibleRef{列上 8:27}\par

祂身穿肉身住在地上，与人的同住是由义人之间的结合与和谐来实现的，他们建造并修筑一座新的圣殿。因为义人就是土地，仍被土地所包围；而且与上主的伟大相比，他们也是土地。因此，有福的伯多禄毫不犹豫地说：\BibleQuote{你们也就像活石，被建成一座属灵的房屋，一座圣洁的殿宇，藉着耶稣基督献上蒙天主悦纳的属灵祭献。}\BibleRef{伯前 2:5}\par

关于那身体，祂以限定的方式将它分别为自己在地上的一块圣地，祂说：\BibleQuote{你们拆毁这座圣殿，三天之内我要把它重建起来。犹太人便说：这座圣殿建筑了四十六年，你能在三天之内把它重建起来吗？但祂所说的圣殿，是指祂的身体。}\BibleRef{若 2:19}\par

% structure: p0219 source=h2 target=section reason=relative-heading-rank; base=h2

\section{四、选自最信奉基督的君王图书馆编号 2431 的手稿。——同上，第 1336 a 页。选自极神圣有福的亚历山大城司铎、斯卓玛提人克莱孟所著\WorkTitle{论天主眷顾}之书。}

什么是天主？正如主所说：\Quote{天主是神。}如今，灵按其确切意义是实体，非物质的，且不受限定的。所谓非物质，即不由身体构成，或其存在不依循宽度、长度与深度。所谓不受限定，即没有处所，全然在万有之中，且完整地临在于每一事物，其自身保持同一。\par

% structure: p0221 source=h2 target=section reason=relative-heading-rank; base=h2

\section{五、选自同一手稿。——同上，1335 D。}

\OriginalTerm{本性}{\GreekText{Φύσις}}之所以得名，源于\OriginalTerm{被生出}{\GreekText{τὸ} \GreekText{πεφυκέναι}}。第一种实体是一切自存之物，例如石头被称为实体。第二种是能生长之实体，如植物生长衰败。第三种是有生命有感觉之实体，如动物、马。第四种是有生命、有感觉、有理性之实体，如人。因此我们每人被造时包含这一切，拥有非物质之灵魂和理智，即天主的肖像。\par

% structure: p0223 source=h2 target=section reason=relative-heading-rank; base=h2

\section{六、选自大马士革的若望\WorkTitle{平行集}卷二，第 307 页。}

对无受情之天主的敬畏不受扰乱。因为人畏惧的不是天主，而是远离天主。畏惧远离的人，畏惧陷入邪恶，也畏惧邪恶本身。畏惧堕落者希望自己是不死不衰的。\par

% structure: p0225 source=h2 target=section reason=relative-heading-rank; base=h2

\section{七、同上，第 341 页。}

应立一条法律，以惩治那些胆敢不敬、不相称地观看神圣与属天事物的人，使他们受到眼盲的刑罚。\par

% structure: p0227 source=h2 target=section reason=relative-heading-rank; base=h2

\section{八、同上，第 657 页。}

总的来说，基督徒是喜爱独处、安静、宁静与和平的。\par

% structure: p0229 source=h2 target=section reason=relative-heading-rank; base=h2

\section{九、选自弗兰西斯·泽菲鲁斯以拉丁文出版的\WorkTitle{梅瑟五书连珠集}，第 146 页。}

那被称作\OriginalTerm{四字圣名}{Tetragrammaton}的奥秘之名，唯有进入至圣所的人方能以此为护，其发音为\OriginalTerm{耶和华}{Jehovah}，意为\Quote{那存在且将要存在者。}祭台南面的灯台象征七颗行星，它们在我们看来似乎围绕中天旋转，两侧各伸出三根枝子；因为太阳也如灯盏，凭天主的上智平衡于众行星之中，以其光照亮上方和下方。祭台另一面放着陈设饼的桌子，因为从那天际吹来维持生命与滋养的和风。\par

% structure: p0231 source=h2 target=section reason=relative-heading-rank; base=h2

\section{十、选自 J. A. 克拉默\WorkTitle{希腊教父新约连珠集}1840 年牛津版卷三。}

论 \BibleBook{宗徒}{大事录} 七章二十四节。奥秘家们说，梅瑟仅凭言语就击杀了那埃及人；正如后来在宗徒大事录中记载，\Person{伯多禄}也凭言语击杀了那些私自扣留田地价银而撒谎的人。\BibleRef{宗 7:24}\par

% structure: p0233 source=h2 target=section reason=relative-heading-rank; base=h2

\section{II. 同上，卷四，第 291 页。}

论 \BibleBook{罗马}{书} 八章三十八节。\BibleQuote{或是生命，指我们现世的生命，}而\BibleQuote{死亡}指迫害者攻击所致的死亡，\BibleQuote{或是天使、掌权者、异能者}指叛离的诸灵。\BibleRef{罗 8:38}\par

% structure: p0235 source=h2 target=section reason=relative-heading-rank; base=h2

\section{十二、第 369 页，第十章第三节。}

他们既不明白也不遵行法律的要求，他们所想象的，也以为法律要求如此。他们不相信法律是预言，只相信字面话；他们出于恐惧而跟从它，而非出于他们的心志与信德。\par

% structure: p0237 source=h2 target=section reason=relative-heading-rank; base=h2

\section{十三、卷六，第 385 页。}

论 \BibleBook{格林多}{后书} 五章十六节。\BibleQuote{即使我们曾按血肉认识基督。}\BibleRef{格后 5:16}\par

他说，至此以后没有人再按血肉生活。因为那不是生命，而是死亡。基督为显明此事，也停止了按血肉生活。如何？不是脱下身体！断乎不是！因为祂将带着这身体作为自己的，作为审判万有的主宰而来。而是除去肉体上的感受，如饥饿、口渴、睡眠和疲倦。因为如今祂拥有一个不能受苦、不能受伤的身体。\par

正如在我们身上，\Quote{按肉体}是指处于罪恶之中，而脱离罪恶便是\Quote{不按肉体}；同样，在基督身上，\Quote{按肉体}是指祂顺从自然情感，而不顺从这些情感便是\Quote{不按肉体}。他却说：\Quote{但正如祂被释放，我们也同样被释放。}他说，再不要顺从肉体的影响。以上是克莱孟在《纲要》第四卷中的论述。\par

% structure: p0241 source=h2 target=section reason=relative-heading-rank; base=h2

\section{十四、出自同一著作，第391页。}

论《格后》6:11：\Quote{我们的心扩大了。}\BibleRef{格后 6:11}\par

因为正如热力通常会使事物膨胀，爱情也是如此。爱情是一种温暖的事物。好像他要说：我不仅用口爱你们，也用心灵，并且将你们全都容纳在内。因此他说：\Quote{你们在我们心中并不狭窄，因为渴望本身扩展了灵魂。}\Quote{我们的心扩大了}为了教导你们一切；\Quote{但你们在自己的肠肚里却狭窄了}，也就是说，在对天主的爱中，你们应当爱我。\par

以上是克莱孟在《纲要》第四卷中的论述。\par

% structure: p0245 source=h2 target=section reason=relative-heading-rank; base=h2

\section{十五、出自第三卷第286页。}

希1:1：\Quote{古时多次并多种方式。}\BibleRef{希 1:1}\par

由于主是全能者的宗徒，被派遣到希伯来人中，保禄出于谦逊，没有自称为希伯来人的宗徒，出于对主的敬畏，以及他是外邦人的宣讲者和宗徒，因此他在本职工作之外又写了《希伯来人书》。\par

% structure: p0248 source=h2 target=section reason=relative-heading-rank; base=h2

\section{十六、出自同一著作。}

同一作品包含出自《导师》第一卷第六章的一段文字。这段文字以\Quote{因为血被发现是}开始，到\Quote{情感的有力魅力}结束。然而，部分内容被省略了。存在相当多的异文；尽管在克拉默著作中发现的这段文字在米涅版本中被完整印出，其理由据称是与克莱孟文本有相当大的差异，但差异并不足以使克拉默版本中的这段文字的翻译具有任何特殊趣味或价值。\par

我们注意到以下异文：——\par

\OriginalTerm{\GreekText{γινεται}}{\GreekText{γινεται}}，动词被省略处，我们插入了\Emph{是}：存在阻碍，等等。\par

\OriginalTerm{\GreekText{σύριγγας}}{\GreekText{σύριγγας}}，管，取代\OriginalTerm{\GreekText{σήραγγας}}{\GreekText{σήραγγας}}（空洞），乳房的空洞。\par

\OriginalTerm{\GreekText{γειτνιαζουσῶν}}{\GreekText{γειτνιαζουσῶν}}，取代\OriginalTerm{\GreekText{γειτνιουσῶν}}{\GreekText{γειτνιουσῶν}}，邻近的（动脉）。\par

\OriginalTerm{\GreekText{ἐπιλήψει}}{\GreekText{ἐπιλήψει}}，取代\OriginalTerm{\GreekText{ἐμπεριλήψει}}{\GreekText{ἐμπεριλήψει}}，中断（如此类的）。\par

\OriginalTerm{\GreekText{ἀποκλήρωσις}}{\GreekText{ἀποκλήρωσις}} 正如文本中出现的那样，注释中指定的校订\OriginalTerm{\GreekText{ἀπολήρησις}}{\GreekText{ἀπολήρησις}} 已被采纳。\par

此处省略的\OriginalTerm{\GreekText{ἥτις} \GreekText{ἐστί}}{\GreekText{ἥτις} \GreekText{ἐστί}}，即\Quote{因恩宠而甘甜}，已补充。\par

% structure: p0257 source=h2 target=section reason=relative-heading-rank; base=h2

\section{第142页。}

\OriginalTerm{\GreekText{γάλα}}{\GreekText{γάλα}}，奶，取代\OriginalTerm{\GreekText{μάννα}}{\GreekText{μάννα}}，玛纳，（那食物）玛纳。\par

% structure: p0259 source=h2 target=section reason=relative-heading-rank; base=h2

\section{第149页。}

\OriginalTerm{\GreekText{χρὴ} \GreekText{δὲ} \GreekText{κατανοῆσαι} \GreekText{τὴν} \GreekText{φύσιν}}{\GreekText{χρὴ} \GreekText{δὲ} \GreekText{κατανοῆσαι} \GreekText{τὴν} \GreekText{φύσιν}}（但必须考虑本性），取代\OriginalTerm{\GreekText{οὐ} \GreekText{κατανενοηκότες}}{\GreekText{οὐ} \GreekText{κατανενοηκότες}}，\OriginalTerm{\GreekText{τ}}{\GreekText{τ}}．\OriginalTerm{\GreekText{φ}}{\GreekText{φ}}．，由于缺乏对本性的考虑。\par

\OriginalTerm{\GreekText{κατακλειομένῳ}}{\GreekText{κατακλειομένῳ}}，与食物一致，取代\OriginalTerm{\GreekText{κατακλειομένω}}{\GreekText{κατακλειομένω}}，与热量一致（封闭在内）。\par

\OriginalTerm{\GreekText{γίνεται}}{\GreekText{γίνεται}} 取代\OriginalTerm{\GreekText{γὰρ}}{\GreekText{γὰρ}}（未翻译），（血）是（准备）用于奶。\par

% structure: p0263 source=h2 target=section reason=relative-heading-rank; base=h2

\section{第144页。}

\OriginalTerm{\GreekText{τοίνυν} \GreekText{τὸν} \GreekText{λόγον}}{\GreekText{τοίνυν} \GreekText{τὸν} \GreekText{λόγον}} 被补充，\OriginalTerm{\GreekText{εἰκότως}}{\GreekText{εἰκότως}} 在从句中被省略，保禄使用恰当的比喻性语言。\par

% structure: p0265 source=h2 target=section reason=relative-heading-rank; base=h2

\section{第145页。}

在\OriginalTerm{\GreekText{ἀλλὰ} \GreekText{τὸ} \GreekText{έν} \GreekText{αὐτῇ}}{\GreekText{ἀλλὰ} \GreekText{τὸ} \GreekText{έν} \GreekText{αὐτῇ}} 之前补充了\OriginalTerm{\GreekText{πλὴ} \GreekText{ν}}{\GreekText{πλὴ} \GreekText{ν}}，而其中的血等被省略。\par

% structure: p0267 source=h2 target=section reason=relative-heading-rank; base=h2

\section{第146页。}

\Quote{因为第欧根尼·阿波罗尼亚特斯主张}被省略。\par

\OriginalTerm{\GreekText{πάντη}}{\GreekText{πάντη}}，译作\Quote{在各方面}，与前一句相连。\par

% structure: p0270 source=h2 target=section reason=relative-heading-rank; base=h2

\section{第147页。}

\OriginalTerm{\GreekText{ὅτι} \GreekText{τοίνυν}}{\GreekText{ὅτι} \GreekText{τοίνυν}}，取代\OriginalTerm{\GreekText{Ὡς} \GreekText{δ᾿}}{\GreekText{Ὡς} \GreekText{δ᾿}}。而（奶被产生）。\par

\OriginalTerm{\GreekText{τηνικαῦτα}}{\GreekText{τηνικαῦτα}} 取代\OriginalTerm{\GreekText{τηνικάδε}}{\GreekText{τηνικάδε}} 在从句中，\Quote{而草和牧场多汁湿润}，未翻译。\par

\OriginalTerm{\GreekText{προειρημένῳ}}{\GreekText{προειρημένῳ}}，上文提到的（奶），省略。\par

\OriginalTerm{\GreekText{τρυφῆς}}{\GreekText{τρυφῆς}} 取代\OriginalTerm{\GreekText{τροφῆς}}{\GreekText{τροφῆς}}，（甘甜的）营养。\par

在\OriginalTerm{\GreekText{γλυκεῖ}}{\GreekText{γλυκεῖ}}，甘甜（酒）之前省略了\OriginalTerm{\GreekText{τῷ}}{\GreekText{τῷ}}，以及\OriginalTerm{\GreekText{καθάπερ}}{\GreekText{καθάπερ}}，\Quote{正如，当受苦时}。\par

\OriginalTerm{\GreekText{τὸ} \GreekText{λιπαρόν}}{\GreekText{τὸ} \GreekText{λιπαρόν}} 取代\OriginalTerm{\GreekText{τῷ} \GreekText{λιπαρῷ}}{\GreekText{τῷ} \GreekText{λιπαρῷ}}，而\OriginalTerm{\GreekText{ἀριδήλως}}{\GreekText{ἀριδήλως}} 取代\OriginalTerm{\GreekText{ἀριδήλου}}{\GreekText{ἀριδήλου}}，在句子中：\Quote{此外，许多人将奶脂（称为黄油）用于灯，明显地，}等等。\par

注意。\par

［勒努里断定《诠释》并非译自《纲要》，但凯（第473页）认为此论断依据不足。亦见凯的学术注释（第5页）。］\par

\SourceRecord{Translated by William Wilson.；Ante-Nicene Fathers；Vol. 2.；Edited by Alexander Roberts, James Donaldson, and A. Cleveland Coxe.；Buffalo, NY: Christian Literature Publishing Co.,；1885.；<http://www.newadvent.org/fathers/0211.htm>.}
